% Options for packages loaded elsewhere
\PassOptionsToPackage{unicode}{hyperref}
\PassOptionsToPackage{hyphens}{url}
\documentclass[
  a4paper,
]{article}
\usepackage{xcolor}
\usepackage[margin=25mm]{geometry}
\usepackage{amsmath,amssymb}
\setcounter{secnumdepth}{-\maxdimen} % remove section numbering
\usepackage{iftex}
\ifPDFTeX
  \usepackage[T1]{fontenc}
  \usepackage[utf8]{inputenc}
  \usepackage{textcomp} % provide euro and other symbols
\else % if luatex or xetex
  \usepackage{unicode-math} % this also loads fontspec
  \defaultfontfeatures{Scale=MatchLowercase}
  \defaultfontfeatures[\rmfamily]{Ligatures=TeX,Scale=1}
\fi
\usepackage{lmodern}
\ifPDFTeX\else
  % xetex/luatex font selection
\fi
% Use upquote if available, for straight quotes in verbatim environments
\IfFileExists{upquote.sty}{\usepackage{upquote}}{}
\IfFileExists{microtype.sty}{% use microtype if available
  \usepackage[]{microtype}
  \UseMicrotypeSet[protrusion]{basicmath} % disable protrusion for tt fonts
}{}
\makeatletter
\@ifundefined{KOMAClassName}{% if non-KOMA class
  \IfFileExists{parskip.sty}{%
    \usepackage{parskip}
  }{% else
    \setlength{\parindent}{0pt}
    \setlength{\parskip}{6pt plus 2pt minus 1pt}}
}{% if KOMA class
  \KOMAoptions{parskip=half}}
\makeatother
\usepackage{color}
\usepackage{fancyvrb}
\newcommand{\VerbBar}{|}
\newcommand{\VERB}{\Verb[commandchars=\\\{\}]}
\DefineVerbatimEnvironment{Highlighting}{Verbatim}{commandchars=\\\{\}}
% Add ',fontsize=\small' for more characters per line
\newenvironment{Shaded}{}{}
\newcommand{\AlertTok}[1]{\textcolor[rgb]{1.00,0.00,0.00}{\textbf{#1}}}
\newcommand{\AnnotationTok}[1]{\textcolor[rgb]{0.38,0.63,0.69}{\textbf{\textit{#1}}}}
\newcommand{\AttributeTok}[1]{\textcolor[rgb]{0.49,0.56,0.16}{#1}}
\newcommand{\BaseNTok}[1]{\textcolor[rgb]{0.25,0.63,0.44}{#1}}
\newcommand{\BuiltInTok}[1]{\textcolor[rgb]{0.00,0.50,0.00}{#1}}
\newcommand{\CharTok}[1]{\textcolor[rgb]{0.25,0.44,0.63}{#1}}
\newcommand{\CommentTok}[1]{\textcolor[rgb]{0.38,0.63,0.69}{\textit{#1}}}
\newcommand{\CommentVarTok}[1]{\textcolor[rgb]{0.38,0.63,0.69}{\textbf{\textit{#1}}}}
\newcommand{\ConstantTok}[1]{\textcolor[rgb]{0.53,0.00,0.00}{#1}}
\newcommand{\ControlFlowTok}[1]{\textcolor[rgb]{0.00,0.44,0.13}{\textbf{#1}}}
\newcommand{\DataTypeTok}[1]{\textcolor[rgb]{0.56,0.13,0.00}{#1}}
\newcommand{\DecValTok}[1]{\textcolor[rgb]{0.25,0.63,0.44}{#1}}
\newcommand{\DocumentationTok}[1]{\textcolor[rgb]{0.73,0.13,0.13}{\textit{#1}}}
\newcommand{\ErrorTok}[1]{\textcolor[rgb]{1.00,0.00,0.00}{\textbf{#1}}}
\newcommand{\ExtensionTok}[1]{#1}
\newcommand{\FloatTok}[1]{\textcolor[rgb]{0.25,0.63,0.44}{#1}}
\newcommand{\FunctionTok}[1]{\textcolor[rgb]{0.02,0.16,0.49}{#1}}
\newcommand{\ImportTok}[1]{\textcolor[rgb]{0.00,0.50,0.00}{\textbf{#1}}}
\newcommand{\InformationTok}[1]{\textcolor[rgb]{0.38,0.63,0.69}{\textbf{\textit{#1}}}}
\newcommand{\KeywordTok}[1]{\textcolor[rgb]{0.00,0.44,0.13}{\textbf{#1}}}
\newcommand{\NormalTok}[1]{#1}
\newcommand{\OperatorTok}[1]{\textcolor[rgb]{0.40,0.40,0.40}{#1}}
\newcommand{\OtherTok}[1]{\textcolor[rgb]{0.00,0.44,0.13}{#1}}
\newcommand{\PreprocessorTok}[1]{\textcolor[rgb]{0.74,0.48,0.00}{#1}}
\newcommand{\RegionMarkerTok}[1]{#1}
\newcommand{\SpecialCharTok}[1]{\textcolor[rgb]{0.25,0.44,0.63}{#1}}
\newcommand{\SpecialStringTok}[1]{\textcolor[rgb]{0.73,0.40,0.53}{#1}}
\newcommand{\StringTok}[1]{\textcolor[rgb]{0.25,0.44,0.63}{#1}}
\newcommand{\VariableTok}[1]{\textcolor[rgb]{0.10,0.09,0.49}{#1}}
\newcommand{\VerbatimStringTok}[1]{\textcolor[rgb]{0.25,0.44,0.63}{#1}}
\newcommand{\WarningTok}[1]{\textcolor[rgb]{0.38,0.63,0.69}{\textbf{\textit{#1}}}}
\setlength{\emergencystretch}{3em} % prevent overfull lines
\providecommand{\tightlist}{%
  \setlength{\itemsep}{0pt}\setlength{\parskip}{0pt}}
\usepackage{bookmark}
\IfFileExists{xurl.sty}{\usepackage{xurl}}{} % add URL line breaks if available
\urlstyle{same}
\hypersetup{
  hidelinks,
  pdfcreator={LaTeX via pandoc}}

\author{}
\date{}

\begin{document}

\section{Anonymous Equal-Amplitude Composite-Wave Model: Basic Axiom
System
v8}\label{anonymous-equal-amplitude-composite-wave-model-basic-axiom-system-v8}

\textbf{Date:} 2026-07-22\\
\textbf{Author:} Noriaki Kihara\\
\textbf{Status:} Definition paper\\
\textbf{Version DOI:} 10.5281/zenodo.21495422 \textbf{Concept DOI:}
10.5281/zenodo.21315735

Changes in v8 (reorganization of layer 0): (1) Axiom 0 (component
anonymity) and Axiom 0+ (formulation anonymity) are merged into a single
Axiom 0 (anonymity) as Clauses 1 and 2 (no change to the text). (2)
Scale anonymity is incorporated into layer 0 as Axiom 0.5 (first
introduced in Paper 5; Concept DOI: 10.5281/zenodo.21486233). (3) The
former Axiom 0.5 (dimension-wise central projection readout) is
renumbered to Axiom 0.6 (no change to the text). (4) References to the
geometric theoretical background (expository note; Concept DOI:
10.5281/zenodo.21495305) are added to Axiom 0.5 and Axiom 1.

\begin{center}\rule{0.5\linewidth}{0.5pt}\end{center}

\section{1. First Principles}\label{first-principles}

\subsection{Axiom 0: Anonymity}\label{axiom-0-anonymity}

Classification: Axiom

This axiom consists of two clauses (Axiom 0 and Axiom 0+ of v7 and
earlier, merged as clauses; no change to the text).

\subsubsection{Clause 1: Component
Anonymity}\label{clause-1-component-anonymity}

Basic components are not assigned individual names, privileged axes,
privileged signs, or privileged types.

The index \texttt{n} is a convenient label and is not an intrinsic name
of the component.

Forbidden:

\begin{Shaded}
\begin{Highlighting}[]
\NormalTok{making only one specific component the negative{-}sign axis}
\NormalTok{making only one specific component the time axis}
\NormalTok{making only one specific component real{-}valued}
\NormalTok{making only one specific component complex{-}valued}
\NormalTok{giving an intrinsic name to only one specific component}
\end{Highlighting}
\end{Shaded}

\begin{center}\rule{0.5\linewidth}{0.5pt}\end{center}

\subsubsection{Clause 2: Formulation
Anonymity}\label{clause-2-formulation-anonymity}

At the first-principle layer, theory names, formulation names, and
readout-logic names are not assigned intrinsic names.

Equations, projections, or readout rules at the first-principle layer
must not be changed on the basis of physical names, object names,
interaction names, or existing theory classifications.

Forbidden:

\begin{Shaded}
\begin{Highlighting}[]
\NormalTok{adopting a circular{-}motion equation because it is called gravity}
\NormalTok{adopting a hyperbolic equation because it is called Coulomb force}
\NormalTok{erasing signs because it is called mass}
\NormalTok{preserving signs because it is called charge}
\NormalTok{moving a component to a time axis because it is called energy}
\NormalTok{using a different closure equation because the target is a two{-}body system}
\NormalTok{using a different closure equation because the target is a three{-}body system}
\NormalTok{using a special readout rule because the target is an ABC system}
\NormalTok{choosing a right{-}hand{-}side representation to obtain a desired physical name}
\end{Highlighting}
\end{Shaded}

Allowed:

\begin{Shaded}
\begin{Highlighting}[]
\NormalTok{different readout operations applied to the same closure condition}
\NormalTok{different symmetry breaking applied to the same closure condition}
\NormalTok{different gauge stability applied to the same closure condition}
\NormalTok{different sign preservation applied to the same closure condition}
\NormalTok{different sign erasure applied to the same closure condition}
\NormalTok{different phase{-}branch selection applied to the same closure condition}
\NormalTok{different closure target sets applied to the same closure condition}
\end{Highlighting}
\end{Shaded}

Order:

\begin{Shaded}
\begin{Highlighting}[]
\NormalTok{place the same all{-}positive{-}sign zero closure}
\NormalTok{define the readout operation first}
\NormalTok{define symmetry breaking first}
\NormalTok{define sign preservation or sign erasure first}
\NormalTok{confirm gauge stability}
\NormalTok{assign working names after readout}
\end{Highlighting}
\end{Shaded}

Target set:

\begin{Shaded}
\begin{Highlighting}[]
\NormalTok{Q(P)=\textbackslash{}sum\_n p\_n\^{}2}
\end{Highlighting}
\end{Shaded}

\begin{Shaded}
\begin{Highlighting}[]
\NormalTok{Q(A)=0}
\end{Highlighting}
\end{Shaded}

\begin{Shaded}
\begin{Highlighting}[]
\NormalTok{Q(A+B)=0}
\end{Highlighting}
\end{Shaded}

\begin{Shaded}
\begin{Highlighting}[]
\NormalTok{Q(A+B+C)=0}
\end{Highlighting}
\end{Shaded}

Differences in target set are allowed.

Changing the definition of \texttt{Q} is not allowed.

\begin{center}\rule{0.5\linewidth}{0.5pt}\end{center}

\subsection{Axiom 0.5: Scale Anonymity}\label{axiom-0.5-scale-anonymity}

Classification: Axiom

The relations, update rules, and readouts of the system must reference
no absolute scale.

Every construction must be covariant under

\begin{Shaded}
\begin{Highlighting}[]
\NormalTok{Z\textbackslash{}to\textbackslash{}lambda Z,\textbackslash{}qquad\textbackslash{}lambda\textbackslash{}in\textbackslash{}mathbb\{C\}\^{}\textbackslash{}ast}
\end{Highlighting}
\end{Shaded}

That is:

\begin{Shaded}
\begin{Highlighting}[]
\NormalTok{relations (closure conditions) are homogeneous}
\NormalTok{the generator is zeroth{-}order homogeneous, K(λZ)=K(Z)}
\NormalTok{the state map is first{-}order homogeneous}
\NormalTok{readouts are ratio quantities only}
\end{Highlighting}
\end{Shaded}

No means exists to measure an absolute scale from inside the system.

This axiom is an unprovable formation rule of the same rank as Axiom 0
(anonymity). The two are clauses of a single idea: no names, no scale
--- only relations and ratios can be read.

The compact part \texttt{\textbar{}λ\textbar{}=1} of the covariance
group \texttt{C*} is the common-phase invariance of the zero closure;
the unreadability of phase (Axiom 5) and the unreadability of scale are
the argument part and the modulus part of one and the same group.

The zero on the right-hand side of Axiom 1 is the unique choice
compatible with this axiom within the family

\begin{Shaded}
\begin{Highlighting}[]
\NormalTok{\textbackslash{}sum\_n x\_n\^{}2=c}
\end{Highlighting}
\end{Shaded}

The first introduction and the discussion are given in Paper 5. Concept
DOI: 10.5281/zenodo.21486233

Theoretical background: writing \texttt{x\_n=q\_n+ip\_n}, the imaginary
part of Axiom 1 is the condition that the generator of scale
transformations (the dilatation) \texttt{Σq\_np\_n} vanishes, so the
identical elimination of the scale degree of freedom is built into the
closure. The geometric organization including this observation is given
in the expository note. Concept DOI: 10.5281/zenodo.21495305

\begin{center}\rule{0.5\linewidth}{0.5pt}\end{center}

\subsection{Axiom 0.6: Dimension-Wise Central Projection
Readout}\label{axiom-0.6-dimension-wise-central-projection-readout}

Classification: Auxiliary geometric axiom (Axiom 0.5 of v7 and earlier;
renumbered upon placing scale anonymity at 0.5; no change to the text)

No curvature correction is introduced for readout along a single axis.

Curvature-correction candidates are introduced only when two or more
independent directions form an area, volume, closed path, or transport
mismatch.

\begin{center}\rule{0.5\linewidth}{0.5pt}\end{center}

\subsection{Axiom 1: All-Positive-Sign Zero
Closure}\label{axiom-1-all-positive-sign-zero-closure}

Classification: Axiom

\begin{Shaded}
\begin{Highlighting}[]
\NormalTok{\textbackslash{}sum\_\{n=1\}\^{}\{N\}x\_n\^{}2=0}
\end{Highlighting}
\end{Shaded}

Expanded form:

\begin{Shaded}
\begin{Highlighting}[]
\NormalTok{x\_1\^{}2+x\_2\^{}2+\textbackslash{}cdots+x\_N\^{}2=0}
\end{Highlighting}
\end{Shaded}

Coefficient signs:

\begin{Shaded}
\begin{Highlighting}[]
\NormalTok{all terms +}
\end{Highlighting}
\end{Shaded}

Not adopted:

\begin{Shaded}
\begin{Highlighting}[]
\NormalTok{\textbackslash{}sum\_n |x\_n|\^{}2=0}
\end{Highlighting}
\end{Shaded}

Term used for closure:

\begin{Shaded}
\begin{Highlighting}[]
\NormalTok{x\_n\^{}2}
\end{Highlighting}
\end{Shaded}

Term not used for closure:

\begin{Shaded}
\begin{Highlighting}[]
\NormalTok{x\_n\textbackslash{}bar\{x\}\_n}
\end{Highlighting}
\end{Shaded}

Applicable states:

\begin{Shaded}
\begin{Highlighting}[]
\NormalTok{stationary closed state}
\NormalTok{stable closed state}
\end{Highlighting}
\end{Shaded}

Quasi-stationary states:

\begin{Shaded}
\begin{Highlighting}[]
\NormalTok{immediately after external influence}
\NormalTok{closure recovery process}
\NormalTok{metastable state}
\end{Highlighting}
\end{Shaded}

Temporary nonzero closure residual in a quasi-stationary state:

\begin{Shaded}
\begin{Highlighting}[]
\NormalTok{\textbackslash{}sum\_n x\_n\^{}2\textbackslash{}ne0}
\end{Highlighting}
\end{Shaded}

After stationarization:

\begin{Shaded}
\begin{Highlighting}[]
\NormalTok{\textbackslash{}sum\_n x\_n\^{}2=0}
\end{Highlighting}
\end{Shaded}

Theoretical background: the solution set of this closure is the complex
isotropic cone, and nontrivial solutions require imaginary parts. Under
the projectivization of Axiom 0.5, the state space becomes a projective
quadric (a compact Kähler manifold), and compactness makes discrete
spectra and finite-dimensional state spaces follow from the geometry.
The expository identification by known mathematics is given in the
expository note. Concept DOI: 10.5281/zenodo.21495305

\begin{center}\rule{0.5\linewidth}{0.5pt}\end{center}

\subsection{Axiom 2: Nontrivial
Existence}\label{axiom-2-nontrivial-existence}

Classification: Axiom

\begin{Shaded}
\begin{Highlighting}[]
\NormalTok{\textbackslash{}exists n,\textbackslash{}quad x\_n\textbackslash{}neq0}
\end{Highlighting}
\end{Shaded}

\begin{center}\rule{0.5\linewidth}{0.5pt}\end{center}

\subsection{Axiom 3: Projection-Axis Degeneration and
Imaginary-Direction
Readout}\label{axiom-3-projection-axis-degeneration-and-imaginary-direction-readout}

Classification: Axiom

When an intrinsic coordinate system is constructed by a projection
method, the projection-axis direction is not distinguishable as a
direction for the residents of that intrinsic coordinate system.

However, if a value in the projection-axis direction remains inside the
closure condition, that value is not erased.

When the compensating quantity in the projection-axis direction is read
as \texttt{z}, it is treated inside the intrinsic coordinate system as

\begin{Shaded}
\begin{Highlighting}[]
\NormalTok{iz}
\end{Highlighting}
\end{Shaded}

This component is unobservable as a first-order direction.

When squared inside the closure expression,

\begin{Shaded}
\begin{Highlighting}[]
\NormalTok{(iz)\^{}2={-}z\^{}2}
\end{Highlighting}
\end{Shaded}

it appears with a reversed sign.

Therefore, an unobservable compensating quantity in the projection-axis
direction appears inside the intrinsic coordinate system as a
sign-reversed compensating quantity through squared readout.

\begin{center}\rule{0.5\linewidth}{0.5pt}\end{center}

\section{2. Consequences from the First
Principles}\label{consequences-from-the-first-principles}

\subsection{Consequence 1: Insufficiency of Real Positive-Definite
Components}\label{consequence-1-insufficiency-of-real-positive-definite-components}

Classification: Derived consequence

\begin{Shaded}
\begin{Highlighting}[]
\NormalTok{x\_n\textbackslash{}in\textbackslash{}mathbb\{R\}}
\NormalTok{\textbackslash{}Rightarrow}
\NormalTok{x\_n\^{}2\textbackslash{}ge0}
\end{Highlighting}
\end{Shaded}

\begin{Shaded}
\begin{Highlighting}[]
\NormalTok{\textbackslash{}sum\_n x\_n\^{}2=0}
\NormalTok{\textbackslash{}Rightarrow}
\NormalTok{x\_n=0\textbackslash{}quad\textbackslash{}forall n}
\end{Highlighting}
\end{Shaded}

\begin{center}\rule{0.5\linewidth}{0.5pt}\end{center}

\subsection{Consequence 2: Internal Generation of Sign
Reversal}\label{consequence-2-internal-generation-of-sign-reversal}

Classification: Derived consequence

\begin{Shaded}
\begin{Highlighting}[]
\NormalTok{x\_1=A,\textbackslash{}qquad x\_2=iA}
\end{Highlighting}
\end{Shaded}

\begin{Shaded}
\begin{Highlighting}[]
\NormalTok{x\_1\^{}2+x\_2\^{}2=A\^{}2+(iA)\^{}2=A\^{}2{-}A\^{}2=0}
\end{Highlighting}
\end{Shaded}

\begin{Shaded}
\begin{Highlighting}[]
\NormalTok{i\^{}2={-}1}
\end{Highlighting}
\end{Shaded}

\begin{center}\rule{0.5\linewidth}{0.5pt}\end{center}

\subsection{Consequence 3: Necessity of Phase
Algebra}\label{consequence-3-necessity-of-phase-algebra}

Classification: Derived consequence

\begin{Shaded}
\begin{Highlighting}[]
\NormalTok{x\_n=r\_n e\^{}\{i\textbackslash{}phi\_n\}}
\end{Highlighting}
\end{Shaded}

\begin{Shaded}
\begin{Highlighting}[]
\NormalTok{x\_n\^{}2=r\_n\^{}2e\^{}\{i2\textbackslash{}phi\_n\}}
\end{Highlighting}
\end{Shaded}

Closure condition:

\begin{Shaded}
\begin{Highlighting}[]
\NormalTok{\textbackslash{}sum\_\{n=1\}\^{}\{N\}r\_n\^{}2e\^{}\{i2\textbackslash{}phi\_n\}=0}
\end{Highlighting}
\end{Shaded}

Equal-amplitude condition:

\begin{Shaded}
\begin{Highlighting}[]
\NormalTok{r\_n=r}
\end{Highlighting}
\end{Shaded}

\begin{Shaded}
\begin{Highlighting}[]
\NormalTok{\textbackslash{}sum\_\{n=1\}\^{}\{N\}e\^{}\{i2\textbackslash{}phi\_n\}=0}
\end{Highlighting}
\end{Shaded}

\begin{center}\rule{0.5\linewidth}{0.5pt}\end{center}

\subsection{Consequence 4: Type
Anonymity}\label{consequence-4-type-anonymity}

Classification: Derived consequence

\begin{Shaded}
\begin{Highlighting}[]
\NormalTok{x\_n\textbackslash{}in\textbackslash{}mathbb\{C\}\textbackslash{}quad\textbackslash{}forall n}
\end{Highlighting}
\end{Shaded}

or

\begin{Shaded}
\begin{Highlighting}[]
\NormalTok{all components are assigned the same phase{-}algebra type}
\end{Highlighting}
\end{Shaded}

\begin{center}\rule{0.5\linewidth}{0.5pt}\end{center}

\subsection{Consequence 5: Ninety-Degree Phase
Difference}\label{consequence-5-ninety-degree-phase-difference}

Classification: Derived consequence

\begin{Shaded}
\begin{Highlighting}[]
\NormalTok{A\^{}2+(iA)\^{}2=0}
\end{Highlighting}
\end{Shaded}

\begin{center}\rule{0.5\linewidth}{0.5pt}\end{center}

\section{3. Anonymous Equal-Amplitude Composite
Wave}\label{anonymous-equal-amplitude-composite-wave}

\subsection{Axiom 4: Anonymous Equal-Amplitude Composite
Wave}\label{axiom-4-anonymous-equal-amplitude-composite-wave}

Classification: Axiom

\begin{Shaded}
\begin{Highlighting}[]
\NormalTok{\textbackslash{}Psi\_N=\textbackslash{}frac\{A\_0\}\{N\}\textbackslash{}sum\_\{k=1\}\^{}\{N\}e\^{}\{i\textbackslash{}theta\_k\}}
\end{Highlighting}
\end{Shaded}

\begin{Shaded}
\begin{Highlighting}[]
\NormalTok{|\textbackslash{}psi\_k|=\textbackslash{}frac\{A\_0\}\{N\}}
\end{Highlighting}
\end{Shaded}

When all phases are aligned:

\begin{Shaded}
\begin{Highlighting}[]
\NormalTok{\textbackslash{}max|\textbackslash{}Psi\_N|=A\_0}
\end{Highlighting}
\end{Shaded}

\begin{center}\rule{0.5\linewidth}{0.5pt}\end{center}

\subsection{Axiom 5: Unreadability of Individual
Phase}\label{axiom-5-unreadability-of-individual-phase}

Classification: Axiom

\begin{Shaded}
\begin{Highlighting}[]
\NormalTok{\textbackslash{}theta\_k\textbackslash{}rightarrow\textbackslash{}theta\_k+\textbackslash{}alpha}
\end{Highlighting}
\end{Shaded}

\begin{center}\rule{0.5\linewidth}{0.5pt}\end{center}

\subsection{Axiom 6: Phase-Difference
Readout}\label{axiom-6-phase-difference-readout}

Classification: Axiom

\begin{Shaded}
\begin{Highlighting}[]
\NormalTok{\textbackslash{}Delta\_\{ij\}=\textbackslash{}theta\_i{-}\textbackslash{}theta\_j}
\end{Highlighting}
\end{Shaded}

Number of readable pair channels:

\begin{Shaded}
\begin{Highlighting}[]
\NormalTok{\textbackslash{}binom\{N\}\{2\}=\textbackslash{}frac\{N(N{-}1)\}\{2\}}
\end{Highlighting}
\end{Shaded}

\begin{center}\rule{0.5\linewidth}{0.5pt}\end{center}

\subsection{Axiom 7: System
Identification}\label{axiom-7-system-identification}

Classification: Axiom

\begin{Shaded}
\begin{Highlighting}[]
\NormalTok{(N,A\_0,\textbackslash{}\{\textbackslash{}Delta\_\{ij\}\textbackslash{}\})}
\end{Highlighting}
\end{Shaded}

\begin{center}\rule{0.5\linewidth}{0.5pt}\end{center}

\subsection{Axiom 8: Composite Phase
Projection}\label{axiom-8-composite-phase-projection}

Classification: Axiom

\begin{Shaded}
\begin{Highlighting}[]
\NormalTok{S\_N=\textbackslash{}sum\_k e\^{}\{i\textbackslash{}theta\_k\}}
\end{Highlighting}
\end{Shaded}

\begin{Shaded}
\begin{Highlighting}[]
\NormalTok{\textbackslash{}Psi\_N=\textbackslash{}frac\{A\_0\}\{N\}S\_N}
\end{Highlighting}
\end{Shaded}

\begin{Shaded}
\begin{Highlighting}[]
\NormalTok{\textbackslash{}Theta\_N=\textbackslash{}arg\textbackslash{}Psi\_N}
\end{Highlighting}
\end{Shaded}

\begin{center}\rule{0.5\linewidth}{0.5pt}\end{center}

\subsection{Axiom 9: No Requirement of Inversion
Symmetry}\label{axiom-9-no-requirement-of-inversion-symmetry}

Classification: Axiom

Inversion symmetry is not required.

\begin{center}\rule{0.5\linewidth}{0.5pt}\end{center}

\section{4. External Readout}\label{external-readout}

\subsection{Axiom 10: External-Axis
Response}\label{axiom-10-external-axis-response}

Classification: Axiom candidate

\begin{Shaded}
\begin{Highlighting}[]
\NormalTok{M\_\textbackslash{}alpha=e\^{}\{i\textbackslash{}alpha\}}
\end{Highlighting}
\end{Shaded}

\begin{Shaded}
\begin{Highlighting}[]
\NormalTok{Y(\textbackslash{}alpha)=\textbackslash{}frac\{A\_0\}\{N\}\textbackslash{}sum\_\{k=1\}\^{}\{N\}e\^{}\{i(\textbackslash{}theta\_k{-}\textbackslash{}alpha)\}}
\NormalTok{=e\^{}\{{-}i\textbackslash{}alpha\}\textbackslash{}Psi\_N}
\end{Highlighting}
\end{Shaded}

Continuous response:

\begin{Shaded}
\begin{Highlighting}[]
\NormalTok{S\_N(\textbackslash{}alpha)=\textbackslash{}operatorname\{Re\}\textbackslash{}left[\textbackslash{}frac\{1\}\{N\}\textbackslash{}sum\_\{k=1\}\^{}\{N\}e\^{}\{i(\textbackslash{}theta\_k{-}\textbackslash{}alpha)\}\textbackslash{}right]}
\end{Highlighting}
\end{Shaded}

Binary response:

\begin{Shaded}
\begin{Highlighting}[]
\NormalTok{S\_N(\textbackslash{}alpha)=\textbackslash{}operatorname\{sgn\}\textbackslash{}operatorname\{Re\}\textbackslash{}left[\textbackslash{}frac\{1\}\{N\}\textbackslash{}sum\_\{k=1\}\^{}\{N\}e\^{}\{i(\textbackslash{}theta\_k{-}\textbackslash{}alpha)\}\textbackslash{}right]}
\end{Highlighting}
\end{Shaded}

\begin{center}\rule{0.5\linewidth}{0.5pt}\end{center}

\subsection{Axiom 11: External Readout
Phase}\label{axiom-11-external-readout-phase}

Classification: Axiom candidate

\begin{Shaded}
\begin{Highlighting}[]
\NormalTok{\textbackslash{}Delta\textbackslash{}phi\_\textbackslash{}mu=\textbackslash{}phi\_\textbackslash{}mu\^{}\{(\textbackslash{}Psi)\}{-}\textbackslash{}phi\_\textbackslash{}mu\^{}\{(M)\}}
\end{Highlighting}
\end{Shaded}

\begin{Shaded}
\begin{Highlighting}[]
\NormalTok{\textbackslash{}mu=x,y,z,t}
\end{Highlighting}
\end{Shaded}

\begin{center}\rule{0.5\linewidth}{0.5pt}\end{center}

\subsection{Axiom 12: Position-Phase
Readout}\label{axiom-12-position-phase-readout}

Classification: Axiom candidate

\begin{Shaded}
\begin{Highlighting}[]
\NormalTok{x,y,z}
\end{Highlighting}
\end{Shaded}

\begin{center}\rule{0.5\linewidth}{0.5pt}\end{center}

\subsection{Axiom 13: Time-Phase
Readout}\label{axiom-13-time-phase-readout}

Classification: Axiom candidate

\begin{Shaded}
\begin{Highlighting}[]
\NormalTok{\textbackslash{}phi\_t\textbackslash{}equiv\textbackslash{}omega t\_\{\textbackslash{}mathrm\{read\}\}\textbackslash{}pmod\{2\textbackslash{}pi\}}
\end{Highlighting}
\end{Shaded}

\begin{center}\rule{0.5\linewidth}{0.5pt}\end{center}

\subsection{Axiom 14: Covering Phase}\label{axiom-14-covering-phase}

Classification: Axiom candidate

\begin{Shaded}
\begin{Highlighting}[]
\NormalTok{2pi{-}periodic readout}
\NormalTok{4pi{-}periodic readout}
\NormalTok{winding{-}number{-}difference readout}
\end{Highlighting}
\end{Shaded}

\begin{center}\rule{0.5\linewidth}{0.5pt}\end{center}

\section{5. Observational Phase
Discreteness}\label{observational-phase-discreteness}

\subsection{Axiom 15: Observational Phase
Discreteness}\label{axiom-15-observational-phase-discreteness}

Classification: Axiom candidate

\begin{Shaded}
\begin{Highlighting}[]
\NormalTok{\textbackslash{}sum\_\{\textbackslash{}mathrm\{cycle\}\}\textbackslash{}Delta\textbackslash{}phi=2\textbackslash{}pi m,\textbackslash{}qquad m\textbackslash{}in\textbackslash{}mathbb\{Z\}}
\end{Highlighting}
\end{Shaded}

\begin{Shaded}
\begin{Highlighting}[]
\NormalTok{W=\textbackslash{}frac\{1\}\{2\textbackslash{}pi\}\textbackslash{}oint d\textbackslash{}phi,\textbackslash{}qquad W\textbackslash{}in\textbackslash{}mathbb\{Z\}}
\end{Highlighting}
\end{Shaded}

Readout-cell form:

\begin{Shaded}
\begin{Highlighting}[]
\NormalTok{\textbackslash{}phi\_\{\textbackslash{}mathrm\{obs\}\}=m\textbackslash{}epsilon\_\textbackslash{}phi,\textbackslash{}qquad m\textbackslash{}in\textbackslash{}mathbb\{Z\}}
\end{Highlighting}
\end{Shaded}

Phase-area candidate:

\begin{Shaded}
\begin{Highlighting}[]
\NormalTok{Q\_\textbackslash{}phi=\textbackslash{}sum\_\{i\textless{}j\}\textbackslash{}sin\^{}2\textbackslash{}frac\{\textbackslash{}Delta\_\{ij\}\}\{2\}}
\end{Highlighting}
\end{Shaded}

\begin{center}\rule{0.5\linewidth}{0.5pt}\end{center}

\section{6. Waveform and Localization}\label{waveform-and-localization}

\subsection{Working Axiom A: State
Waveform}\label{working-axiom-a-state-waveform}

Classification: Working hypothesis

\begin{Shaded}
\begin{Highlighting}[]
\NormalTok{A(\textbackslash{}theta)=\textbackslash{}sum\_n A\_n\textbackslash{}cos(n\textbackslash{}theta+\textbackslash{}phi\_n)}
\end{Highlighting}
\end{Shaded}

\begin{Shaded}
\begin{Highlighting}[]
\NormalTok{Z(\textbackslash{}theta)=\textbackslash{}sum\_n r\_n e\^{}\{i(n\textbackslash{}theta+\textbackslash{}phi\_n)\}}
\end{Highlighting}
\end{Shaded}

\begin{center}\rule{0.5\linewidth}{0.5pt}\end{center}

\subsection{Definition A1: Normalized Equal-Amplitude Odd-Harmonic
Localized
Wave}\label{definition-a1-normalized-equal-amplitude-odd-harmonic-localized-wave}

Classification: Definition

\begin{Shaded}
\begin{Highlighting}[]
\NormalTok{N\_h=2K{-}1}
\end{Highlighting}
\end{Shaded}

\begin{Shaded}
\begin{Highlighting}[]
\NormalTok{S\_\{N\_h\}(u)=\textbackslash{}frac\{1\}\{K\}\textbackslash{}sum\_\{m=0\}\^{}\{K{-}1\}\textbackslash{}cos((2m+1)u)}
\end{Highlighting}
\end{Shaded}

\begin{Shaded}
\begin{Highlighting}[]
\NormalTok{Z\_\{N\_h\}(u)=\textbackslash{}frac\{1\}\{K\}\textbackslash{}sum\_\{m=0\}\^{}\{K{-}1\}e\^{}\{i(2m+1)u\}}
\end{Highlighting}
\end{Shaded}

\begin{Shaded}
\begin{Highlighting}[]
\NormalTok{\textbackslash{}Psi\_\{N\_h\}(u)=A S\_\{N\_h\}(u)}
\end{Highlighting}
\end{Shaded}

\begin{Shaded}
\begin{Highlighting}[]
\NormalTok{\textbackslash{}Psi\_\{N\_h\}(u)=A Z\_\{N\_h\}(u)}
\end{Highlighting}
\end{Shaded}

\begin{Shaded}
\begin{Highlighting}[]
\NormalTok{\textbackslash{}frac\{A\}\{K\}}
\end{Highlighting}
\end{Shaded}

\begin{center}\rule{0.5\linewidth}{0.5pt}\end{center}

\subsection{Theorem A1: Representative-Amplitude
Preservation}\label{theorem-a1-representative-amplitude-preservation}

Classification: Derived theorem

\begin{Shaded}
\begin{Highlighting}[]
\NormalTok{S\_\{N\_h\}(0)=1}
\end{Highlighting}
\end{Shaded}

\begin{Shaded}
\begin{Highlighting}[]
\NormalTok{\textbackslash{}Psi\_\{N\_h\}(0)=A}
\end{Highlighting}
\end{Shaded}

\begin{center}\rule{0.5\linewidth}{0.5pt}\end{center}

\subsection{Theorem A2: Highest Odd-Harmonic Order and Localization
Width}\label{theorem-a2-highest-odd-harmonic-order-and-localization-width}

Classification: Derived theorem

\begin{Shaded}
\begin{Highlighting}[]
\NormalTok{K=\textbackslash{}frac\{N\_h+1\}\{2\}}
\end{Highlighting}
\end{Shaded}

\begin{Shaded}
\begin{Highlighting}[]
\NormalTok{\textbackslash{}lambda\_\{N\_h\}=\textbackslash{}frac\{\textbackslash{}lambda\_0\}\{N\_h\}}
\end{Highlighting}
\end{Shaded}

\begin{Shaded}
\begin{Highlighting}[]
\NormalTok{\textbackslash{}Delta x\_\{N\_h\}\textbackslash{}propto\textbackslash{}frac\{\textbackslash{}lambda\_0\}\{N\_h\}}
\end{Highlighting}
\end{Shaded}

\begin{Shaded}
\begin{Highlighting}[]
\NormalTok{N\_\{h,\textbackslash{}rm req\}\textbackslash{}sim\textbackslash{}frac\{\textbackslash{}lambda\_0\}\{\textbackslash{}Delta x\}}
\end{Highlighting}
\end{Shaded}

\begin{Shaded}
\begin{Highlighting}[]
\NormalTok{N\_\{h,\textbackslash{}rm req\}\^{}\{(t)\}\textbackslash{}sim\textbackslash{}frac\{T\_0\}\{\textbackslash{}Delta t\}}
\end{Highlighting}
\end{Shaded}

\begin{center}\rule{0.5\linewidth}{0.5pt}\end{center}

\subsection{Consequence A1: Identical Construction of Spatial and
Temporal
Localization}\label{consequence-a1-identical-construction-of-spatial-and-temporal-localization}

Classification: Derived consequence

\begin{Shaded}
\begin{Highlighting}[]
\NormalTok{\textbackslash{}Psi(\textbackslash{}chi,\textbackslash{}tau)=A S\_\{N\_\{h,\textbackslash{}chi\}\}(\textbackslash{}chi{-}\textbackslash{}chi\_0)S\_\{N\_\{h,\textbackslash{}tau\}\}(\textbackslash{}tau{-}\textbackslash{}tau\_0)e\^{}\{i\textbackslash{}phi\}}
\end{Highlighting}
\end{Shaded}

\begin{center}\rule{0.5\linewidth}{0.5pt}\end{center}

\subsection{Working Axiom B:
Localization}\label{working-axiom-b-localization}

Classification: Working hypothesis

\begin{Shaded}
\begin{Highlighting}[]
\NormalTok{S\_\{N\_h\}(u)=\textbackslash{}frac\{1\}\{K\}\textbackslash{}sum\_\{m=0\}\^{}\{K{-}1\}\textbackslash{}cos((2m+1)u)}
\end{Highlighting}
\end{Shaded}

\begin{Shaded}
\begin{Highlighting}[]
\NormalTok{Z\_\{N\_h\}(u)=\textbackslash{}frac\{1\}\{K\}\textbackslash{}sum\_\{m=0\}\^{}\{K{-}1\}e\^{}\{i(2m+1)u\}}
\end{Highlighting}
\end{Shaded}

\begin{center}\rule{0.5\linewidth}{0.5pt}\end{center}

\subsection{Working Axiom C: Coupling
Entrance}\label{working-axiom-c-coupling-entrance}

Classification: Working hypothesis

The entrance to interaction is low-order base-phase alignment.

\begin{center}\rule{0.5\linewidth}{0.5pt}\end{center}

\subsection{Working Axiom D: Observation
Map}\label{working-axiom-d-observation-map}

Classification: Working hypothesis

\begin{Shaded}
\begin{Highlighting}[]
\NormalTok{s\_\textbackslash{}alpha=\textbackslash{}mathrm\{Loc\}\textbackslash{}left[\textbackslash{}int A(\textbackslash{}theta)M\_\textbackslash{}alpha(\textbackslash{}theta)d\textbackslash{}theta\textbackslash{}right]}
\end{Highlighting}
\end{Shaded}

\begin{center}\rule{0.5\linewidth}{0.5pt}\end{center}

\subsection{Working Axiom E: Squared Correlation
Strength}\label{working-axiom-e-squared-correlation-strength}

Classification: Working hypothesis

\begin{Shaded}
\begin{Highlighting}[]
\NormalTok{P(\textbackslash{}alpha)=\textbackslash{}frac\{|C\_\textbackslash{}alpha|\^{}2\}\{\textbackslash{}sum\_\textbackslash{}beta |C\_\textbackslash{}beta|\^{}2\}}
\end{Highlighting}
\end{Shaded}

\begin{center}\rule{0.5\linewidth}{0.5pt}\end{center}

\subsection{Working Axiom F: Shared Low-Order Base
Synchronization}\label{working-axiom-f-shared-low-order-base-synchronization}

Classification: Working hypothesis

\begin{Shaded}
\begin{Highlighting}[]
\NormalTok{\textbackslash{}theta\_A{-}\textbackslash{}theta\_B=\textbackslash{}Delta}
\end{Highlighting}
\end{Shaded}

\begin{center}\rule{0.5\linewidth}{0.5pt}\end{center}

\subsection{Working Axiom G: Synchronized Demodulation
Failure}\label{working-axiom-g-synchronized-demodulation-failure}

Classification: Working hypothesis

\begin{Shaded}
\begin{Highlighting}[]
\NormalTok{V(t)=\textbackslash{}left|\textbackslash{}sum\_n w\_n e\^{}\{i\textbackslash{}epsilon\_n(t)\}\textbackslash{}right|}
\end{Highlighting}
\end{Shaded}

\begin{Shaded}
\begin{Highlighting}[]
\NormalTok{V(t)=\textbackslash{}left|\textbackslash{}sum\_n w\_n e\^{}\{{-}\textbackslash{}frac12\textbackslash{}sigma\_n\^{}2(t)\}\textbackslash{}right|}
\end{Highlighting}
\end{Shaded}

\begin{center}\rule{0.5\linewidth}{0.5pt}\end{center}

\subsection{Working Axiom H: Hierarchy}\label{working-axiom-h-hierarchy}

Classification: Working hypothesis

\begin{Shaded}
\begin{Highlighting}[]
\NormalTok{low{-}order base mode}
\NormalTok{+ high{-}order localized mode}
\end{Highlighting}
\end{Shaded}

\begin{center}\rule{0.5\linewidth}{0.5pt}\end{center}

\section{7. Readout, Memory, and Complex
Representation}\label{readout-memory-and-complex-representation}

\subsection{I/Q Readout}\label{iq-readout}

Classification: Theoretical observation

\begin{Shaded}
\begin{Highlighting}[]
\NormalTok{z=a+ib}
\end{Highlighting}
\end{Shaded}

\begin{Shaded}
\begin{Highlighting}[]
\NormalTok{a=r\textbackslash{}cos\textbackslash{}theta,\textbackslash{}qquad b=r\textbackslash{}sin\textbackslash{}theta}
\end{Highlighting}
\end{Shaded}

\begin{Shaded}
\begin{Highlighting}[]
\NormalTok{u=r\textbackslash{}cos\textbackslash{}theta+r\textbackslash{}sin\textbackslash{}theta}
\NormalTok{=\textbackslash{}sqrt2 r\textbackslash{}cos\textbackslash{}left(\textbackslash{}theta{-}\textbackslash{}frac\{\textbackslash{}pi\}\{4\}\textbackslash{}right)}
\end{Highlighting}
\end{Shaded}

\begin{center}\rule{0.5\linewidth}{0.5pt}\end{center}

\subsection{Harmonic-Information
Readout}\label{harmonic-information-readout}

Classification: Theoretical observation

\begin{Shaded}
\begin{Highlighting}[]
\NormalTok{Z(\textbackslash{}theta)=\textbackslash{}sum\_n r\_n e\^{}\{i(n\textbackslash{}theta+\textbackslash{}phi\_n)\}}
\end{Highlighting}
\end{Shaded}

\begin{Shaded}
\begin{Highlighting}[]
\NormalTok{P(Z)=\textbackslash{}operatorname\{Re\}(Z)+\textbackslash{}operatorname\{Im\}(Z)}
\end{Highlighting}
\end{Shaded}

\begin{Shaded}
\begin{Highlighting}[]
\NormalTok{P(Z)=\textbackslash{}sum\_n \textbackslash{}sqrt2 r\_n}
\NormalTok{\textbackslash{}cos\textbackslash{}left(n\textbackslash{}theta+\textbackslash{}phi\_n{-}\textbackslash{}frac\{\textbackslash{}pi\}\{4\}\textbackslash{}right)}
\end{Highlighting}
\end{Shaded}

\begin{center}\rule{0.5\linewidth}{0.5pt}\end{center}

\subsection{Preservation Side and Readout
Side}\label{preservation-side-and-readout-side}

Classification: Theoretical observation

\begin{Shaded}
\begin{Highlighting}[]
\NormalTok{z\textbackslash{}bar z=a\^{}2+b\^{}2}
\end{Highlighting}
\end{Shaded}

\begin{Shaded}
\begin{Highlighting}[]
\NormalTok{z\^{}2=a\^{}2{-}b\^{}2+2iab}
\end{Highlighting}
\end{Shaded}

\begin{center}\rule{0.5\linewidth}{0.5pt}\end{center}

\subsection{Square Registry}\label{square-registry}

Classification: Working hypothesis

\begin{Shaded}
\begin{Highlighting}[]
\NormalTok{\textbackslash{}mathcal\{A\}\_0=\textbackslash{}sum\_k A\_k\^{}2}
\end{Highlighting}
\end{Shaded}

\begin{Shaded}
\begin{Highlighting}[]
\NormalTok{a\_0=\textbackslash{}sum\_k |A\_k|\^{}2=\textbackslash{}sum\_k A\_k\textbackslash{}bar A\_k}
\end{Highlighting}
\end{Shaded}

\begin{center}\rule{0.5\linewidth}{0.5pt}\end{center}

\section{8. Statistical
Classification}\label{statistical-classification}

\subsection{Pair Phase Difference}\label{pair-phase-difference}

Classification: Definition

\begin{Shaded}
\begin{Highlighting}[]
\NormalTok{\textbackslash{}Delta\_\{ij\}=\textbackslash{}theta\_i{-}\textbackslash{}theta\_j}
\end{Highlighting}
\end{Shaded}

\begin{Shaded}
\begin{Highlighting}[]
\NormalTok{\textbackslash{}binom\{N\}\{2\}=\textbackslash{}frac\{N(N{-}1)\}\{2\}}
\end{Highlighting}
\end{Shaded}

Pair composite:

\begin{Shaded}
\begin{Highlighting}[]
\NormalTok{\textbackslash{}psi\_i+\textbackslash{}psi\_j}
\NormalTok{=2A\textbackslash{}cos\textbackslash{}frac\{\textbackslash{}Delta\_\{ij\}\}\{2\}e\^{}\{i(\textbackslash{}theta\_i+\textbackslash{}theta\_j)/2\}}
\end{Highlighting}
\end{Shaded}

Pair interference strength:

\begin{Shaded}
\begin{Highlighting}[]
\NormalTok{I\_\{ij\}=4A\^{}2\textbackslash{}cos\^{}2\textbackslash{}frac\{\textbackslash{}Delta\_\{ij\}\}\{2\}}
\end{Highlighting}
\end{Shaded}

\begin{center}\rule{0.5\linewidth}{0.5pt}\end{center}

\subsection{Overlap Degree}\label{overlap-degree}

Classification: Definition

\begin{Shaded}
\begin{Highlighting}[]
\NormalTok{R\_N=\textbackslash{}left|\textbackslash{}frac\{A\_0\}\{N\}\textbackslash{}sum\_\{k=1\}\^{}\{N\}e\^{}\{i\textbackslash{}theta\_k\}\textbackslash{}right|}
\end{Highlighting}
\end{Shaded}

\begin{Shaded}
\begin{Highlighting}[]
\NormalTok{\textbackslash{}rho\_N=\textbackslash{}frac\{R\_N\}\{A\_0\}}
\NormalTok{=\textbackslash{}frac\{1\}\{N\}\textbackslash{}left|\textbackslash{}sum\_\{k=1\}\^{}\{N\}e\^{}\{i\textbackslash{}theta\_k\}\textbackslash{}right|}
\end{Highlighting}
\end{Shaded}

\begin{Shaded}
\begin{Highlighting}[]
\NormalTok{0\textbackslash{}le\textbackslash{}rho\_N\textbackslash{}le1}
\end{Highlighting}
\end{Shaded}

\begin{center}\rule{0.5\linewidth}{0.5pt}\end{center}

\subsection{Fermionic Degree}\label{fermionic-degree}

Classification: Definition

\begin{Shaded}
\begin{Highlighting}[]
\NormalTok{F\_N=}
\NormalTok{\textbackslash{}frac\{2\}\{N(N{-}1)\}}
\NormalTok{\textbackslash{}sum\_\{i\textless{}j\}\textbackslash{}sin\^{}2\textbackslash{}frac\{\textbackslash{}Delta\_\{ij\}\}\{2\}}
\end{Highlighting}
\end{Shaded}

\begin{Shaded}
\begin{Highlighting}[]
\NormalTok{B\_N=1{-}F\_N}
\end{Highlighting}
\end{Shaded}

\begin{Shaded}
\begin{Highlighting}[]
\NormalTok{B\_N=}
\NormalTok{\textbackslash{}frac\{2\}\{N(N{-}1)\}}
\NormalTok{\textbackslash{}sum\_\{i\textless{}j\}\textbackslash{}cos\^{}2\textbackslash{}frac\{\textbackslash{}Delta\_\{ij\}\}\{2\}}
\end{Highlighting}
\end{Shaded}

\begin{Shaded}
\begin{Highlighting}[]
\NormalTok{F\_N=\textbackslash{}frac\{N\}\{2(N{-}1)\}(1{-}\textbackslash{}rho\_N\^{}2)}
\end{Highlighting}
\end{Shaded}

\begin{center}\rule{0.5\linewidth}{0.5pt}\end{center}

\section{9. Observation Selection and Curvature
Projection}\label{observation-selection-and-curvature-projection}

\subsection{Axiom 16: Unique Observation Selection by Complete Pairwise
Relational
Waves}\label{axiom-16-unique-observation-selection-by-complete-pairwise-relational-waves}

Classification: observation-connection axiom

This axiom is an independent observation principle not derived from
Axiom 0 or Axiom 1.

For a closed system with component count \texttt{N}, define the
unordered complete pairwise relation set without proper names as

\begin{Shaded}
\begin{Highlighting}[]
\NormalTok{\textbackslash{}mathcal E\_N}
\NormalTok{=\textbackslash{}left\textbackslash{}\{\textbackslash{}\{i,j\textbackslash{}\}\textbackslash{}mid 1\textbackslash{}le i\textless{}j\textbackslash{}le N\textbackslash{}right\textbackslash{}\}}
\end{Highlighting}
\end{Shaded}

For each relation

\begin{Shaded}
\begin{Highlighting}[]
\NormalTok{e=\textbackslash{}\{i,j\textbackslash{}\}\textbackslash{}in\textbackslash{}mathcal E\_N}
\end{Highlighting}
\end{Shaded}

the corresponding

\begin{Shaded}
\begin{Highlighting}[]
\NormalTok{X\_e\textbackslash{}in\textbackslash{}mathbb C}
\end{Highlighting}
\end{Shaded}

is a physical relational wave that constitutes the system, not an
observation-device channel.

The relational wave \texttt{X\_e} is not a quantity derived from
component waves; it is an independent physical state component. The
correspondence map to component waves is not specified by this axiom.

Write the relational-wave state as

\begin{Shaded}
\begin{Highlighting}[]
\NormalTok{X\_\{\textbackslash{}mathcal E\}}
\NormalTok{=\textbackslash{}left(X\_e\textbackslash{}right)\_\{e\textbackslash{}in\textbackslash{}mathcal E\_N\}}
\end{Highlighting}
\end{Shaded}

A permutation of individual names acts as a permutation of relational
waves and does not change the closure equation, the update rule, or the
observation selection.

Define the quadratic form of the relational-wave set as

\begin{Shaded}
\begin{Highlighting}[]
\NormalTok{q\_\{\textbackslash{}mathcal E\}(X)}
\NormalTok{=\textbackslash{}sum\_\{e\textbackslash{}in\textbackslash{}mathcal E\_N\}X\_e\^{}2}
\end{Highlighting}
\end{Shaded}

The closure of the relational waves is written, using the compensation
quantity of Axiom 3, as

\begin{Shaded}
\begin{Highlighting}[]
\NormalTok{\textbackslash{}sum\_\{e\textbackslash{}in\textbackslash{}mathcal E\_N\}X\_e\^{}2+(iR)\^{}2=0}
\end{Highlighting}
\end{Shaded}

The transposed form

\begin{Shaded}
\begin{Highlighting}[]
\NormalTok{q\_\{\textbackslash{}mathcal E\}(X)=R\^{}2}
\end{Highlighting}
\end{Shaded}

may be used.

Let the number of relational waves be

\begin{Shaded}
\begin{Highlighting}[]
\NormalTok{M=\textbackslash{}left|\textbackslash{}mathcal E\_N\textbackslash{}right|}
\end{Highlighting}
\end{Shaded}

and let a real-coefficient generator acting identically on all
relational waves be

\begin{Shaded}
\begin{Highlighting}[]
\NormalTok{K\textbackslash{}in\textbackslash{}mathbb R\^{}\{M\textbackslash{}times M\}}
\end{Highlighting}
\end{Shaded}

The relational wave \texttt{X\_e} remains complex, and \texttt{K} acts
as the same real-coefficient transformation on its real and imaginary
parts.

Write the local linear update of relational waves as

\begin{Shaded}
\begin{Highlighting}[]
\NormalTok{\textbackslash{}dot X=KX}
\end{Highlighting}
\end{Shaded}

If this update preserves the quadratic form identically,

\begin{Shaded}
\begin{Highlighting}[]
\NormalTok{\textbackslash{}frac\{d\}\{d\textbackslash{}tau\}q\_\{\textbackslash{}mathcal E\}(X)}
\NormalTok{=X\^{}\{\textbackslash{}mathsf T\}\textbackslash{}left(K\^{}\{\textbackslash{}mathsf T\}+K\textbackslash{}right)X}
\NormalTok{=0}
\end{Highlighting}
\end{Shaded}

Therefore, the closure-preserving generator of relational waves
satisfies

\begin{Shaded}
\begin{Highlighting}[]
\NormalTok{K\^{}\{\textbackslash{}mathsf T\}={-}K}
\end{Highlighting}
\end{Shaded}

This antisymmetry is not introduced from any particular physical name or
axis name.

A real antisymmetric generator decomposes, without additional proper
names, privileged axes, or external references, into a direct sum of
two-dimensional rotation planes and a null space.

Let the number of rotation planes be

\begin{Shaded}
\begin{Highlighting}[]
\NormalTok{r=\textbackslash{}frac12\textbackslash{}operatorname\{rank\}K}
\end{Highlighting}
\end{Shaded}

and the null-space dimension be

\begin{Shaded}
\begin{Highlighting}[]
\NormalTok{\textbackslash{}dim\textbackslash{}ker K=M{-}2r}
\end{Highlighting}
\end{Shaded}

The adoption conditions for intrinsically observable directions are as
follows.

First, rotation planes become candidates for distinct directions only
when they can be mutually distinguished by the relational quantity of
independent rotation frequencies. Distinguishability is a necessary
condition and does not by itself license adoption as spatial directions.

Second, when a normal direction is read from the plane spanned by two
linearly independent relational directions, the normal may be adopted as
a direction unique up to sign only when the normal-candidate space is
one-dimensional. For a real \texttt{d}-dimensional spatial display, the
dimension of the normal-candidate space is

\begin{Shaded}
\begin{Highlighting}[]
\NormalTok{d{-}2}
\end{Highlighting}
\end{Shaded}

and this condition holds only for

\begin{Shaded}
\begin{Highlighting}[]
\NormalTok{d=3}
\end{Highlighting}
\end{Shaded}

Third, the number of intrinsically observable directions is not counted
from the relational-wave count \texttt{M}, the generator rank, or the
null-space dimension. It is counted as the number of directions
satisfying the uniqueness conditions above.

The minimal example is \texttt{M=3}: when

\begin{Shaded}
\begin{Highlighting}[]
\NormalTok{\textbackslash{}operatorname\{rank\}K=2,}
\NormalTok{\textbackslash{}qquad}
\NormalTok{\textbackslash{}dim\textbackslash{}ker K=1}
\end{Highlighting}
\end{Shaded}

the two directions of one rotation plane and the one normal direction
uniquely determined from that plane, three directions in total, may be
adopted.

When \texttt{M=6} and three rotation planes are distinguished by
independent rotation frequencies, three directions may be adopted.

When the number of distinguishable rotation planes exceeds the range of
spatial display satisfying the second uniqueness condition, the excess
rotation planes are not adopted as spatial directions and are retained
as internal phase modes.

When multiple isomorphic rotation planes remain, or when a
normal-candidate space or null space of dimension two or more remains
and its internal directions cannot be distinguished by relational
quantities alone, no single direction among them may be selected as an
observable spatial direction. Such a residue is retained as an internal
subspace without a unique basis.

The primitive readout of each relational wave is constructed from the
phase difference of Axiom 6,

\begin{Shaded}
\begin{Highlighting}[]
\NormalTok{\textbackslash{}Delta\_\{ij\}=\textbackslash{}theta\_i{-}\textbackslash{}theta\_j}
\end{Highlighting}
\end{Shaded}

and the bilinear relation obtained from the quadratic form of Axiom 1,

\begin{Shaded}
\begin{Highlighting}[]
\NormalTok{q(X)=\textbackslash{}sum\_n X\_n\^{}2}
\end{Highlighting}
\end{Shaded}

\begin{Shaded}
\begin{Highlighting}[]
\NormalTok{B(U,V)}
\NormalTok{=\textbackslash{}frac\{1\}\{2\}\textbackslash{}left(q(U+V){-}q(U){-}q(V)\textbackslash{}right)}
\end{Highlighting}
\end{Shaded}

Readout for three or more bodies is constructed from the set of complete
pairwise relational waves and the invariant decomposition of their
closure-preserving generator.

No independent higher-order readout rule may be added on the grounds of
physical names, object names, or dimension names.

The specific coupling strengths, phase-difference dependence, and update
period of the relational-wave generator are not specified by this axiom.

They are defined as independent relational update rules consistent with
Axiom 0, Axiom 0.5, Axiom 1, and this axiom.

\begin{center}\rule{0.5\linewidth}{0.5pt}\end{center}

\subsection{Working Axiom I: Phase-Difference Sine Relational Update
Rule}\label{working-axiom-i-phase-difference-sine-relational-update-rule}

Classification: working hypothesis

As a concrete form of the relational-wave generator of Axiom 16, define
the endpoint-sharing adjacency

\begin{Shaded}
\begin{Highlighting}[]
\NormalTok{A\_\{ef\}}
\NormalTok{=}
\NormalTok{\textbackslash{}begin\{cases\}}
\NormalTok{1,\& e\textbackslash{}ne f\textbackslash{} \textbackslash{}text\{and\}\textbackslash{} e\textbackslash{}cap f\textbackslash{}ne\textbackslash{}varnothing,\textbackslash{}\textbackslash{}}
\NormalTok{0,\& \textbackslash{}text\{otherwise\}}
\NormalTok{\textbackslash{}end\{cases\}}
\end{Highlighting}
\end{Shaded}

and, from the initial phases \texttt{theta\_e\ =\ arg\ X\_e}, set

\begin{Shaded}
\begin{Highlighting}[]
\NormalTok{\textbackslash{}widetilde K\_\{ef\}}
\NormalTok{=}
\NormalTok{A\_\{ef\}\textbackslash{}sin(\textbackslash{}theta\_f{-}\textbackslash{}theta\_e)}
\end{Highlighting}
\end{Shaded}

It satisfies

\begin{Shaded}
\begin{Highlighting}[]
\NormalTok{\textbackslash{}widetilde K\^{}\{\textbackslash{}mathsf T\}={-}\textbackslash{}widetilde K}
\end{Highlighting}
\end{Shaded}

and thus the antisymmetry condition of Axiom 16.

This coupling rule is not derived from Axiom 0, Axiom 0.5, Axiom 1, or
Axiom 16. It is a working hypothesis placed as one concrete
closure-preserving relational update rule. Note that this generator, as
a function of phase differences only, is zeroth-order homogeneous and
satisfies Axiom 0.5 without adjustment (Paper 5, Theorem 3. Concept DOI:
10.5281/zenodo.21486233).

\begin{center}\rule{0.5\linewidth}{0.5pt}\end{center}

\subsection{Consequence 16.1: Vertex Decomposition and Linear Rank
Bound}\label{consequence-16.1-vertex-decomposition-and-linear-rank-bound}

Classification: derived consequence (under Working Axiom I)

Let \texttt{B} be the unsigned incidence matrix with individuals as rows
and relations as columns, and from its \texttt{k}-th row build

\begin{Shaded}
\begin{Highlighting}[]
\NormalTok{c\_k=(B\_\{ke\}\textbackslash{}cos\textbackslash{}theta\_e)\_\{e\},}
\NormalTok{\textbackslash{}qquad}
\NormalTok{s\_k=(B\_\{ke\}\textbackslash{}sin\textbackslash{}theta\_e)\_\{e\}}
\end{Highlighting}
\end{Shaded}

Then

\begin{Shaded}
\begin{Highlighting}[]
\NormalTok{\textbackslash{}widetilde K=\textbackslash{}sum\_\{k=1\}\^{}\{N\}\textbackslash{}left(c\_k s\_k\^{}\{\textbackslash{}mathsf T\}{-}s\_k c\_k\^{}\{\textbackslash{}mathsf T\}\textbackslash{}right)}
\end{Highlighting}
\end{Shaded}

holds exactly.

Each term is an antisymmetric matrix of rank at most 2, so for all
\texttt{N} and all initial phases,

\begin{Shaded}
\begin{Highlighting}[]
\NormalTok{\textbackslash{}operatorname\{rank\}\textbackslash{}widetilde K}
\NormalTok{\textbackslash{}le}
\NormalTok{2\textbackslash{}min\textbackslash{}!\textbackslash{}left(N,\textbackslash{}left\textbackslash{}lfloor\textbackslash{}frac\{M\}\{2\}\textbackslash{}right\textbackslash{}rfloor\textbackslash{}right)}
\end{Highlighting}
\end{Shaded}

The number of relational waves

\begin{Shaded}
\begin{Highlighting}[]
\NormalTok{M=\textbackslash{}binom\{N\}\{2\}}
\end{Highlighting}
\end{Shaded}

grows quadratically, while the number of rotation modes

\begin{Shaded}
\begin{Highlighting}[]
\NormalTok{r=\textbackslash{}frac12\textbackslash{}operatorname\{rank\}\textbackslash{}widetilde K}
\end{Highlighting}
\end{Shaded}

grows at most linearly.

Proof: Paper 3 of the Dimensional Generation Structure series. Concept
DOI: 10.5281/zenodo.21465898

\begin{center}\rule{0.5\linewidth}{0.5pt}\end{center}

\subsection{Consequence 16.2: Generic-Position Rank
Equality}\label{consequence-16.2-generic-position-rank-equality}

Classification: derived consequence (computer-assisted proof,
\texttt{3\ \textless{}=\ N\ \textless{}=\ 12})

For each \texttt{N} with \texttt{3\ \textless{}=\ N\ \textless{}=\ 12},
for all initial phases except a set of Lebesgue measure zero,

\begin{Shaded}
\begin{Highlighting}[]
\NormalTok{\textbackslash{}operatorname\{rank\}\textbackslash{}widetilde K}
\NormalTok{=}
\NormalTok{2\textbackslash{}min\textbackslash{}!\textbackslash{}left(N,\textbackslash{}left\textbackslash{}lfloor\textbackslash{}frac\{M\}\{2\}\textbackslash{}right\textbackslash{}rfloor\textbackslash{}right)}
\end{Highlighting}
\end{Shaded}

The proof combines exact rational witness configurations via the tan
half-angle parametrization with the fact that the zero set of a
nontrivial real-analytic function has Lebesgue measure zero. Paper 3.
Concept DOI: 10.5281/zenodo.21465898

The generic null-space dimension is then

\begin{Shaded}
\begin{Highlighting}[]
\NormalTok{\textbackslash{}dim\textbackslash{}ker\textbackslash{}widetilde K}
\NormalTok{=}
\NormalTok{M{-}2\textbackslash{}min\textbackslash{}!\textbackslash{}left(N,\textbackslash{}left\textbackslash{}lfloor\textbackslash{}frac\{M\}\{2\}\textbackslash{}right\textbackslash{}rfloor\textbackslash{}right)}
\end{Highlighting}
\end{Shaded}

which for \texttt{N\ \textgreater{}=\ 5} equals

\begin{Shaded}
\begin{Highlighting}[]
\NormalTok{\textbackslash{}frac\{N(N{-}5)\}\{2\}}
\end{Highlighting}
\end{Shaded}

Equality for all \texttt{N} is a new hypothesis and is not included in
this consequence.

\begin{center}\rule{0.5\linewidth}{0.5pt}\end{center}

\subsection{Consequence 16.3: Normal
Uniqueness}\label{consequence-16.3-normal-uniqueness}

Classification: derived consequence

When a normal direction is reconstructed from the plane spanned by two
linearly independent relational directions, without additional
background axes or selection rules, the normal-candidate space of a real
\texttt{d}-dimensional display has dimension

\begin{Shaded}
\begin{Highlighting}[]
\NormalTok{d{-}2}
\end{Highlighting}
\end{Shaded}

and the normal is unique up to sign only for

\begin{Shaded}
\begin{Highlighting}[]
\NormalTok{d=3}
\end{Highlighting}
\end{Shaded}

For \texttt{d\ \textgreater{}=\ 4}, the orthogonal group

\begin{Shaded}
\begin{Highlighting}[]
\NormalTok{O(d{-}2)}
\end{Highlighting}
\end{Shaded}

acting on the normal-candidate space while fixing the plane leaves a
continuous selection freedom, and under anonymity (Axiom 0) no single
direction can be selected. The projector and squared quantities are
unique, but the individual linear directions inside are not.

\subsection{This consequence is the basis of the second adoption
condition of Axiom 16. Proof: Paper 3. Concept DOI:
10.5281/zenodo.21465898}\label{this-consequence-is-the-basis-of-the-second-adoption-condition-of-axiom-16.-proof-paper-3.-concept-doi-10.5281zenodo.21465898}

\subsection{Consequence 16.4: Plane Decomposition of Fixed Generators
and Kinematic
Isomorphism}\label{consequence-16.4-plane-decomposition-of-fixed-generators-and-kinematic-isomorphism}

Classification: derived consequence (under a fixed generator)

A fixed closure-preserving generator decomposes into a direct sum of
two-dimensional rotation planes and a null space. A Cayley-type
orthogonal update acts on each rotation plane as a rotation by the
constant angle

\begin{Shaded}
\begin{Highlighting}[]
\NormalTok{\textbackslash{}theta\_j=2\textbackslash{}arctan(\textbackslash{}gamma\textbackslash{}sigma\_j)}
\end{Highlighting}
\end{Shaded}

and on the null space as the identity. The coordinate of each plane is
therefore a constant-amplitude single-frequency rotation, kinematically
isomorphic to the rotational kinematics of a single two-body relation.

The orientation of each plane is fixed by the generator itself through

\begin{Shaded}
\begin{Highlighting}[]
\NormalTok{Kp\_j=\textbackslash{}sigma\_j q\_j}
\end{Highlighting}
\end{Shaded}

The residual gauge preserving the standard form is \texttt{SO(2)}; on
planes of nonzero amplitude the signed phase progression, the frequency,
and the amplitude are readable, while the absolute phase is not. This
means that the unreadability of Axiom 5 is reproduced also at the level
of the derived plane coordinates. If reflections are allowed
(\texttt{O(2)}), the invariants are the absolute values of progression
and frequency.

The null-space component is static and carries no phase progression, so
progression-based readouts do not hold from the null space (the
dynamical realization of Consequence 16.3).

Frequency nondegeneracy is an independent general-position assumption,
not derivable from the generic rank equality (Consequence 16.2).

Proof: Paper 4 of the Dimensional Generation Structure series. Concept
DOI: 10.5281/zenodo.21468959

\begin{center}\rule{0.5\linewidth}{0.5pt}\end{center}

\subsection{Consequence 16.5: Double Conservation
Register}\label{consequence-16.5-double-conservation-register}

Classification: derived consequence (under a fixed generator)

Both the energy-like quantity and the squared closure decompose
plane-wise and null-space-wise, with each term individually conserved
under the orthogonal update of a fixed generator.

\begin{Shaded}
\begin{Highlighting}[]
\NormalTok{H=\textbackslash{}sum\_\{j\}H\_j+H\_\{\textbackslash{}ker\},}
\NormalTok{\textbackslash{}qquad}
\NormalTok{H\_j=\textbackslash{}lvert p\_j\^{}\textbackslash{}dagger X\textbackslash{}rvert\^{}2+\textbackslash{}lvert q\_j\^{}\textbackslash{}dagger X\textbackslash{}rvert\^{}2}
\end{Highlighting}
\end{Shaded}

\begin{Shaded}
\begin{Highlighting}[]
\NormalTok{X\^{}\{\textbackslash{}mathsf T\}X=\textbackslash{}sum\_\{j\}Q\_j+Q\_\{\textbackslash{}ker\}=R\^{}2,}
\NormalTok{\textbackslash{}qquad}
\NormalTok{Q\_j=(p\_j\^{}\{\textbackslash{}mathsf T\}X)\^{}2+(q\_j\^{}\{\textbackslash{}mathsf T\}X)\^{}2}
\end{Highlighting}
\end{Shaded}

The fixed-generator system thus carries a double conservation register
\texttt{(H\_j,\ Q\_j)} on each plane and on the null space.

Proof: Paper 4. Concept DOI: 10.5281/zenodo.21468959

\begin{center}\rule{0.5\linewidth}{0.5pt}\end{center}

\subsection{Consequence 16.6: No-Go for Spontaneous
Splitting}\label{consequence-16.6-no-go-for-spontaneous-splitting}

Classification: derived consequence (under a fixed generator)

Under a fixed generator, neither the closure quantities \texttt{Q\_j}
nor the energy-like quantities \texttt{H\_j} move between planes, or
between planes and the null space.

Spontaneous splitting of the state accompanied by redistribution among
registers therefore cannot occur under fixed-generator dynamics.
Describing splitting requires a relational update rule beyond the fixed
generator: sequentially reconstructed generators, nonlinear couplings,
or inter-plane exchange.

This consequence provides the starting boundary for any additional
definition of splitting dynamics consistent with Axiom 0, Axiom 0.5,
Axiom 1, and Axiom 16.

Proof: Paper 4. Concept DOI: 10.5281/zenodo.21468959

\begin{center}\rule{0.5\linewidth}{0.5pt}\end{center}

\subsection{Axiom 17: Curvature Reaction Projection Centered on the
Future Phase
Position}\label{axiom-17-curvature-reaction-projection-centered-on-the-future-phase-position}

Classification: dynamical projection axiom

This axiom is an independent dynamical and projection principle not
derived from Axiom 0 or Axiom 1.

Assign an order index

\begin{Shaded}
\begin{Highlighting}[]
\NormalTok{\textbackslash{}tau\textbackslash{}in\textbackslash{}mathbb\{N\}}
\end{Highlighting}
\end{Shaded}

to the state sequence of a closed system.

In the state update

\begin{Shaded}
\begin{Highlighting}[]
\NormalTok{X(\textbackslash{}tau)\textbackslash{}longrightarrow X(\textbackslash{}tau+1)}
\end{Highlighting}
\end{Shaded}

the phase position on the \texttt{tau+1} side is defined as the future
phase position of that state sequence.

When the phase progression on a rotation plane uniquely constructed
without additional names by Axiom 16 has a curvature-radius readout
\texttt{rho\_c} and an angular-velocity readout \texttt{omega}, and can
be expressed as a motion whose relational virtual rotation center is the
future phase position, the center-directed compensation that closes this
virtual rotation --- that is, the reaction to the centrifugal force ---
is mapped and read out as an observable acceleration in the
circumferential direction of progression.

The center-directed compensation is the internal representation, and the
circumferential acceleration is the external readout. This axiom defines
the projection connecting the two as a rule.

The magnitude of the acceleration is

\begin{Shaded}
\begin{Highlighting}[]
\NormalTok{|\textbackslash{}alpha\_\{\textbackslash{}mathrm\{read\}\}|}
\NormalTok{=\textbackslash{}rho\_c|\textbackslash{}omega|\^{}2}
\end{Highlighting}
\end{Shaded}

Let \texttt{t\_hat\_tau} be the unit vector of the circumferential
direction of progression pointing from the current phase position toward
the future phase position. The readout acceleration is

\begin{Shaded}
\begin{Highlighting}[]
\NormalTok{\textbackslash{}boldsymbol\{\textbackslash{}alpha\}\_\{\textbackslash{}mathrm\{read\}\}}
\NormalTok{=\textbackslash{}rho\_c|\textbackslash{}omega|\^{}2\textbackslash{},\textbackslash{}hat\{\textbackslash{}boldsymbol\{t\}\}\_\textbackslash{}tau}
\end{Highlighting}
\end{Shaded}

Here \texttt{rho\_c}, \texttt{omega}, and \texttt{t\_hat\_tau} are
relational readout quantities without physical names in the
first-principles layer. The circumferential direction is not an absolute
axis of a background space; it is determined each time from the relation
among the current phase position, the future phase position, and the
closure radius.

This axiom may not be modified on the grounds of gravity, Coulomb force,
mass, charge, or any other physical name.

The same projection rule applies to every closure mode satisfying the
same conditions.

If the phase branch opposite to the future side, a reversal of the
acceleration direction, or sign retention or sign elimination is
adopted, it must not be selected from physical names; it must be defined
in advance as an independent phase relation, branch selection, symmetry
breaking, or readout operation.

This axiom does not identify the readout acceleration with standard
gravity, the standard Coulomb force, or any known external force.

Physical names are given only after the acceleration map, distance
exponent, sign, closure preservation, and gauge stability have been
confirmed.

\end{document}
